\documentclass[11pt]{article}
\usepackage[T1]{fontenc}
\usepackage{lmodern}
\usepackage{royden-report}
\usepackage{microtype,array,longtable,enumitem,float}
\usepackage[font=small,labelfont=bf]{caption}
\setlength{\intextsep}{8pt}
\setlist{nosep,leftmargin=*}
\newcommand{\ii}{\mathrm{i}}
\newcommand{\D}{\lvert D_s\rvert}
\newcommand{\norm}[1]{\left\lVert #1\right\rVert}
\newcommand{\fig}[3]{\begin{figure}[H]\centering\includegraphics[width=#2\linewidth]{figures/#1.pdf}\caption{#3}\end{figure}}
\hypersetup{pdftitle={Dysthe modelling: numerical evidence and a PINN research framework},pdfauthor={Matas Vai research working report}}
\title{Dysthe modelling: numerical evidence and a PINN research framework}
\author{Research meeting report}
\date{30 October 2026}
\begin{document}
\maketitle
\begin{abstract}
The immediate objective is to reproduce and assess finite-time evolution of several smooth water-wave envelope profiles, then determine whether a neural approximation preserves the same amplitude, phase and derivative-sensitive behaviour. This report combines a new CPU numerical proof of concept for a one-coordinate temporal potential-envelope Dysthe model with previously recorded, separately labelled spatial free-surface PINN results. It identifies the equation, verifies controls, measures refinement and states where interpretation must stop. Neither computation nor a small PINN residual establishes global existence of the continuum PDE.
\end{abstract}
\section{What can be presented now}
\begin{enumerate}
\item \textbf{A reproducible original-paper numerical baseline.} A periodic Fourier solver evolves a localized packet, a modulated wavetrain and a smooth random profile. Independent grid and time-step comparisons support all six primary profile--steepness cases through scaled time $T=16$, with a smaller step needed in one case. The bundle includes the runs, diagnostics, exact-solution controls and a CPU cluster launcher.
\item \textbf{A concrete PINN improvement already measured.} On three spatial validation inlets, an invertible Fourier coefficient scaling reduces maximum relative $H^2$ error from $12$--$26\%$ to $0.22$--$1.47\%$ at the same update budget. These are CPU fits on the displayed interval, with one optimizer seed; the repaired A100 campaign is awaiting results.
\item \textbf{A useful failure diagnosis.} Conservation alone does not certify accuracy. Coarse runs retain action while failing derivative-sensitive grid comparisons near focusing. Some finer-grid stress runs also require smaller time steps. An unresolved computation cannot be interpreted as a singularity.
\end{enumerate}
\begin{center}\small
\begin{tabular}{p{.25\linewidth}p{.30\linewidth}p{.36\linewidth}}\toprule
Question & Evidence available & Conclusion justified \\\midrule
Do several profiles evolve reproducibly? & CPU temporal solver, refinement, controls & Finite-time periodic trajectories on checked intervals \\
Can PINN curvature accuracy improve? & Paired spatial CPU fits & Output conditioning is promising for this architecture \\
Does the continuum solution remain smooth forever? & No matching global theorem or certificate supplied & Open within this project; numerical diagnostics guide analysis \\
Does an FNO generalize faster than a solver? & A test protocol, no new validated FNO result & Requires held-out whole-profile and cost comparisons \\\bottomrule
\end{tabular}
\end{center}
\noindent\textbf{Meeting decision.} Confirm the temporal formulation and first finite-time validation window before extending neural training.
\clearpage
\section{The equation and its physical meaning}\label{sec:model}
Dysthe's 1979 paper derives a fourth-order \emph{asymptotic envelope approximation} for weakly nonlinear, narrowband, deep-water gravity waves \cite{dysthe}. Fourth order refers to the steepness expansion, not to a fourth spatial derivative. Its complex envelope contains amplitude and phase; it does not resolve each carrier-wave crest. The original $A$ is a velocity-potential envelope, whereas the paper's $B$ denotes a surface-elevation envelope. These variables must not be interchanged silently.

With the paper's dimensionless coordinates and $\rho=|A|^2$, the literal printed equation (2.19) is
\begin{align}
2\ii(A_t+\tfrac12A_x)+\tfrac12A_{yy}-\tfrac14A_{xx}-\rho A
={}&-\tfrac{\ii}{8}(6A_{xyy}-A_{xxx})\notag\\
&+\tfrac{3\ii}{2}A(A\overline{A}_x-\overline{A}A_x)
 +\tfrac{\ii}{2}\rho A_x+A(\bar\phi_x-\ii\bar\phi_z).
\label{eq:original}
\end{align}
Here $\bar\phi$ is the induced mean potential, not the complex conjugate of $A$. It is harmonic below the surface and decays at depth. At the retained order, the mean-flow boundary condition is $\bar\phi_z=\rho_x$ at $z=0$; neglect of the higher-order mean-elevation time derivative is part of the asymptotic reduction. For a longitudinal restriction ($\partial_y=0$), harmonic extension gives
\[
\widehat{\bar\phi}(k,z)=\ii\operatorname{sgn}(k)\widehat\rho(k)e^{|k|z},
\qquad \bar\phi_x=-|D_x|\rho,
\]
with zero derivative at the zero Fourier mode. Substitution of \emph{both} boundary derivatives is needed to obtain the local derivative coefficients.

\subsection{A documented coefficient ambiguity}
Literal reduction of \eqref{eq:original} yields $c=1$ in
\begin{equation}
 A_t+\tfrac12A_x+\tfrac{\ii}{8}A_{xx}-\tfrac1{16}A_{xxx}
 +\tfrac{\ii}{2}\rho A+c\rho A_x-\tfrac14A^2\overline A_x
 -\tfrac{\ii}{2}A|D_x|\rho=0.\label{eq:temporal}
\end{equation}
However, the printed amplitude--phase equations in the paper's Section 3 require $c=3/2$. Reversing the sign of the displayed $\ii\rho A_x/2$ on the right of \eqref{eq:original} reconciles them. This is an algebraic observation; no authenticated erratum was verified. Later primary discussion of the potential-envelope conjugate term supports $c=3/2$ with the negative conjugate coefficient \cite{shemer,trulsen}.

\textbf{Report convention.} The primary temporal numerical plots use $c=3/2$, labelled \emph{stability-consistent}. The literal $c=1$ equation remains a separately named comparator. Neither is silently substituted for the approved spatial model. The source audit in the bundle records equations, source pages and the limits of verification.

\subsection{Computational variables and a distinct spatial problem}
Set $A=\epsilon u$, $s=\epsilon(x-t/2)$ and $T=\epsilon^2t$. Then \eqref{eq:temporal} becomes
\begin{equation}
 u_T=-\tfrac{\ii}{8}u_{ss}-\tfrac{\ii}{2}|u|^2u
 +\epsilon\left[\tfrac1{16}u_{sss}-c|u|^2u_s
 +\tfrac14u^2\overline{u_s}+\tfrac{\ii}{2}u\D(|u|^2)\right].
 \label{eq:scaled}
\end{equation}
This is a nonlinear, nonlocal temporal equation on a periodic interval of length $L=40$. Its data are initial \emph{spatial profiles} $u(s,0)$. It is a one-coordinate restriction, not a simulation of the full two-horizontal-coordinate original system.
\clearpage
The existing repository and the PINN evidence use the different spatial free-surface model \cite{fedele}
\begin{equation}
 u_\xi=-\ii u_{\tau\tau}-\ii|u|^2u
 -\epsilon\left[8|u|^2u_\tau+2u^2\overline{u_\tau}
                  +2\ii u|D_\tau|(|u|^2)\right].\label{eq:spatial}
\end{equation}
Here $\xi$ is propagation distance and $\tau$ is retarded time; its input is a wavemaker-type envelope trace. Different coefficients and evolution variables encode a different formulation. Equation \eqref{eq:spatial} is the active contract \texttt{water-wave-dysthe-spatial-v1}. New temporal files live in a separately versioned proof-of-concept directory; no active scope check was relaxed.

\section{What later behaviour means}\label{sec:later}
There is no universal final value that a Dysthe solution must attain. In the conservative periodic model, a wave need not settle to a steady profile. Relevant questions concern peak amplification, modulation growth, drift, phase, spectral broadening, repeated focusing and sensitivity to initial data. A visually localized pulse is not automatically a soliton, and a repeating amplitude pattern is not automatically recurrence of the complex field.

For \eqref{eq:scaled}, define the nonlinear time scale using the initial peak $a_*=\max_s|u(s,0)|$:
\begin{equation}
 T_{\mathrm{nl}}=\frac{2}{a_*^2},\qquad
 n_{\mathrm{nl}}=T/T_{\mathrm{nl}}.
\end{equation}
The cubic phase coefficient is $1/2$. For the selected profiles, $T=16$ is about $5.6$--$5.8$ such nonlinear times. This is a later finite-time window, not an asymptotic $T\to\infty$ study. In physical time, $T=\epsilon^2\omega_0t_{\mathrm{phys}}$; a carrier frequency is required before reporting seconds. Temporal $T$ must not be equated with spatial $\xi$.

\begin{definition}[Recorded observables]
For $k_m=2\pi m/L$ and $\widehat u_m=L^{-1}\int_0^L u(s)e^{-\ii k_ms}\,\rad s$, let
\[
 M=L\sum_m|\widehat u_m|^2,\qquad
 \norm{u}_{H^j}^2=L\sum_m(1+k_m^2)^j|\widehat u_m|^2.
\]
The spectral centroid and bandwidth are
\[\mu_k=\frac{\sum_m k_m|\widehat u_m|^2}{\sum_m|\widehat u_m|^2},\qquad\sigma_k^2=\frac{\sum_m(k_m-\mu_k)^2|\widehat u_m|^2}{\sum_m|\widehat u_m|^2}.\]
\end{definition}
We record $\max|u|$, $M$, $H^1$, $H^2$, $\mu_k$, $\sigma_k$, high-mode power and boundary action. The proxies $\epsilon\max|u|$ and $\epsilon\sigma_k$ indicate changing steepness and bandwidth. No universal cutoff is imposed as a validated physical criterion. A well-resolved envelope computation may still leave the narrowband, weak-steepness regime in which its derivation is credible.

\textbf{Two different uses of $H^2$.} Growth of $\norm{u(T)}_{H^2}$ is a property of a trajectory. The prediction error
\begin{equation}
 E_j(T)=\frac{\norm{v(T)-u_{\rm ref}(T)}_{H^j}}{\norm{u_{\rm ref}(T)}_{H^j}}
 \label{eq:error}
\end{equation}
is a property of an approximation relative to a reference. Large $E_2$ does not show a PDE singularity; bounded computed $H^2$ does not prove global smoothness. Neither norm is evaluated after phase or translation alignment unless explicitly stated.

\clearpage
\section{A runnable numerical proof of concept}\label{sec:numerics}
\subsection{Discretization and acceptance rules}
The new solver uses Fourier differentiation and an integrating-factor fourth-order Runge--Kutta method. The linear multiplier in \eqref{eq:scaled} is
\[
 \lambda(k)=\ii k^2/8-\ii\epsilon k^3/16.
\]
Reference arithmetic is complex128. Exact linear evolution does not remove the nonlinear time-step restriction. The strict mask $|m|<N/4$ prevents cubic aliasing into retained modes. The density is formed before its $|D|$ multiplier; final cubic terms are projected, retaining the nonlocal mean flow.

Each accepted trajectory is compared separately with $N/2$ points at the same time step and with the same $N$ at twice the time step. At saved times separated by $0.2$, both comparisons must satisfy
\[
\max E_0\le10^{-5},\quad \max E_1\le10^{-4},\quad
\max E_2\le10^{-3}.
\]
Maximum action drift must be at most $10^{-6}$ and action in the upper fifth of retained Fourier modes at most $10^{-8}$. These are declared empirical acceptance rules, not rigorous error enclosures. The principal $T=16$ audit uses $N=4096$, $h=0.0025$, with comparisons $N=2048$ and $h=0.005$. Smaller grids and all failed outcomes are retained. Time sampling may miss an unsampled extreme.

\subsection{Three inputs, two steepness parameters, four equations}
The common box is $[-20,20)$. The packet is a periodized Gaussian with peak $0.85$, width $3$, centre $-2$, quadratic chirp $0.08$ and global phase $0.3$. The wavetrain is
\[
 0.8\,[1+0.05\cos(qs)+0.025\ii\sin(qs)]e^{-0.2\ii},
 \qquad q=2\pi\cdot5/40.
\]
The random profile is a fixed, smooth complex Fourier polynomial with modes $-5,\ldots,5$, Gaussian-decaying random coefficients and peak normalized to $0.85$. Coefficients are saved. Initial data are projected through mode $12$, which is part of their definition.
\fig{temporal-initial-profiles}{.83}{Inputs for the temporal study. The same profiles are used at $\epsilon=0.05$ and $0.1$. One random realization establishes a reproducible example, not ensemble statistics.}
Cubic NLS suppresses the entire $\epsilon$ bracket in \eqref{eq:scaled}. A mean-flow-only control retains just its last term, following the purpose of Dysthe's simplified Section 4 model. Literal-print and stability-consistent Dysthe retain all terms. Controls change the equation and are labelled accordingly; they are not alternative numerical solvers for identical physics.
\clearpage
\subsection{Observed evolution on checked intervals}
\fig{temporal-evolution}{1}{Temporal potential-envelope dynamics for $c=3/2$, $\epsilon=0.05$. All three trajectories pass the stated independent comparisons through $T=16$. The packet focuses and moves; modulation grows on the wavetrain; the smooth random profile develops a strong localized peak. Display arrays are Fourier-resampled to 256 points; diagnostics use the full reference grid. Grayscale intensity is envelope magnitude, not free-surface elevation.}
The observed localization is meaningful as a finite-time prediction of this envelope model. It is not evidence of wave breaking, a permanent soliton, or a singularity. The small box uses periodic interactions; packet domain checks below only support the earlier interval $T\le8$. A physical free-surface reconstruction would require the appropriate potential-to-elevation relation and carrier parameters.
\clearpage
\subsection{Peak growth and derivative growth quantify different behaviour}
\fig{temporal-observables}{1}{Full-grid diagnostics for the three $\epsilon=0.05$ stability-consistent trajectories. Peaks and $H^2$ norms increase without implying a loss of regularity. The lower panels track changes in formal steepness and bandwidth; neither is a certified criterion for agreement with the full water-wave equations.}
\begin{center}\small
\begin{tabular}{lrrrr}\toprule
Profile & Peak amplification & Max. $H^2/H^2_0$ & First sampled peak $T$ & Max. grid $E_2$ \\\midrule
Packet & 2.593 & 7.284 & 9.8 & $1.24\cdot10^{-9}$ \\
Wavetrain & 2.781 & 7.479 & 15.2 & $1.04\cdot10^{-8}$ \\
Random & 3.026 & 6.574 & 11.8 & $1.11\cdot10^{-6}$ \\\bottomrule
\end{tabular}
\end{center}
The corresponding maximum time-step $E_2$ values are $3.51\cdot10^{-7}$, $2.87\cdot10^{-7}$ and $5.92\cdot10^{-6}$. Maximum relative action drift is below $1.3\cdot10^{-11}$ in all three. These observations support a resolved focusing episode for these particular periodic profiles. They neither bound all higher derivatives nor cover arbitrary data.

\textbf{Domain check.} Doubling the packet period from $40$ to $80$, doubling the grid, and keeping the physical Fourier cutoff matched gives maximum common-interval $H^1$ difference about $3.8\cdot10^{-4}$ at $\epsilon=0.05$ and $7.6\cdot10^{-4}$ at $0.1$, through $T=8$. Both temporal variants pass the declared $1\%$ domain rule. This is a two-box comparison, not a whole-line limit; it does not establish domain independence through $T=16$.
\clearpage
\subsection{Refinement is essential even when conservation looks excellent}
\fig{temporal-refinement}{1}{Grid comparison and action drift for $\epsilon=0.05$. The first two columns cover $T\le8$; the final column covers $T\le16$. Dashed horizontal lines show the respective acceptance thresholds. Conservation and grid agreement test different properties.}
A fine-grid calculation can expose a time-step restriction that a coarser grid did not reveal. Exact treatment of the linear dispersion does not make the nonlinear integrating-factor scheme unconditionally stable. The initial high-resolution audit passes 21 of 24 model--profile--steepness combinations; the three failures are the $\epsilon=0.1$ wavetrain in the literal, stability-consistent and mean-flow-only models. All failures and blocked continuations are saved. A separate smaller-step audit resolves the primary $c=3/2$ failure: $h=0.0003125$, compared with $0.000625$, reduces the apparent maximum $H^2/H^2_0$ from $600.6$ to $28.42$, with grid $E_2=4.92\cdot10^{-9}$ and step $E_2=6.95\cdot10^{-10}$. All six primary profile--steepness cases thus qualify through $T=16$, using this finer step for one case. The two other high-steepness control failures remain unresolved. No earlier failure is overwritten.
\clearpage
\subsection{Dysthe corrections become visible beyond cubic NLS}
\fig{temporal-model-differences}{.80}{Complex-field $L^2$ differences from the $c=3/2$ temporal Dysthe solution at $\epsilon=0.05$. Every displayed model pair passes its own $T=16$ numerical audit. Curves are computed from portable display arrays; the full-grid endpoint comparisons in the text are authoritative. No phase or translation alignment is performed.}
At $T=16$, the full-grid complex $L^2$ difference between the stability-consistent and NLS equations, normalized by the NLS solution, is approximately $0.910$ for the packet, $1.163$ for the wavetrain and $0.892$ for the random profile. The figure uses the stability-consistent solution in the denominator, so it is a distinct normalization. Even with similar amplitude ranges, accumulated phase and drift can make complex-field discrepancies large. These are \emph{model differences}, not neural approximation errors.

The literal and stability-consistent variants also separate substantially: their endpoint relative differences, normalized by the stability-consistent solution, are about $0.217$, $0.514$ and $0.297$. Accordingly, selecting a formulation is consequential for later dynamics. Agreement with a instability growth rate about a zero-envelope-wavenumber background alone cannot resolve the printed coefficient ambiguity.

\subsection{Exact controls before learned models}
For a plane wave $u=a\exp[\ii(qs-\Omega T)]$, the mean-flow derivative vanishes and
\begin{equation}
 \Omega=-q^2/8+\epsilon q^3/16+a^2/2
       +\epsilon(c+1/4)a^2q.\label{eq:plane}
\end{equation}
A nonzero $q$ detects the difference between $c=1$ and $c=3/2$. The bundled tests also check the linear propagator, Fourier mean-flow symbol, cubic dealiasing, zero-$\epsilon$ agreement, time convergence and a cubic-NLS sech solution. The latter is a solver control in the NLS limit, not a claimed Dysthe soliton.

For the zero-carrier-envelope background, the exact retained-model instability exponent obeys
\begin{equation}
 \gamma^2=\frac{q^2}{8}\left[a^2(1-\epsilon|q|)-\frac{q^2}{8}\right]
                  -\frac{\epsilon^2a^4q^2}{16}.\label{eq:mi}
\end{equation}
Inside the unstable band ($\gamma^2>0$), its real sideband frequency is $\epsilon(q^3/16+ca^2q)$. The two $c$ values have identical growth but different drift. Dysthe's formal first-order steepness correction omits the last term in \eqref{eq:mi}; implementation checks should use the exact eigenvalue of the equation actually integrated.

\clearpage
\section{What the existing PINN results establish}\label{sec:pinn}
\subsection{A slab PINN fits short consecutive intervals}
The recorded experiment solves the spatial equation \eqref{eq:spatial}. A separate neural network is fitted for each input profile and each slab of propagation distance. For $\delta=\xi-\xi_j\in[0,0.5]$, the trial field has the form
\begin{equation}
 v_j(\delta,\tau)=e^{\delta L}P_Ku_j(\tau)
  +\delta\sum_{|m|\le K}w_m c_{\theta_j,m}(\delta)e^{\ii k_m\tau},
 \qquad w_m=\frac1{1+k_m^2}.\label{eq:ansatz}
\end{equation}
Here $L=-\ii\partial_{\tau\tau}$, $P_K$ is Fourier projection and the distance-only network returns real and imaginary Fourier coefficients. The projected input is enforced exactly at each slab entrance. The computed endpoint supplies the next input; there is no reset to reference truth. The network is trained on the PDE residual throughout each slab. This is a numerical fit over that slab, not an extrapolation beyond its fitted distance.

The baseline uses $w_m=1$. With all retained $w_m>0$, the change in \eqref{eq:ansatz} is an invertible parameterization on the retained modes; it does not remove high modes, modify the equation or change evaluation gates. It suppresses large initial changes in derivative-sensitive coefficients during optimization. The mechanism is a plausible conditioning explanation, not a proof of an optimizer convergence rate.

Eight slabs cover $\xi\in[0,4]$. Both versions use $K=48$, width 64, depth 3, 400 Adam updates per slab, 64 residual-distance slices and a 512-point spatial residual grid. A finer independent 1024-point residual audit is retained. Checkpoints are selected by the full PDE residual; numerical labels are used for evaluation, not for training these fits. The reference grid is 1024 with time step $0.00125$.
\fig{spatial-pinn-error-growth}{1}{Spatial free-surface equation, three validation inlets, optimizer seed 11. Maximum relative $H^2$ error is taken over all 81 saved stations through $\xi=4$. The 5\% line is one component of the existing acceptance gate. These are fitted-distance results, not forecasts or new A100 outcomes.}
\clearpage
\subsection{The improvement is substantial, but bounded in scope}
\begin{center}\small
\begin{tabular}{lrrr}\toprule
Validation inlet & Baseline max. $E_2$ & Conditioned max. $E_2$ & Reduction factor \\\midrule
Packet & 26.112\% & 1.466\% & 17.8 \\
Sideband & 12.108\% & 0.655\% & 18.5 \\
Random & 15.850\% & 0.216\% & 73.5 \\\bottomrule
\end{tabular}
\end{center}
All three conditioned fits pass the original sampled gate: $E_0\le1\%$, $E_1,E_2\le5\%$ and peak error at most $5\%$. All three unconditioned fits fail. These are validation profiles used during method development, not an untouched test set. Equal update budgets give recorded training times of approximately 14.1--14.4 seconds unconditioned and 15.3--15.7 seconds conditioned. One CPU timing per fit does not establish speedup over a Fourier solver; it only shows the modest recorded overhead of this change.
\fig{spatial-pinn-profiles-xi4}{.76}{Endpoint profiles at $\xi=4$, including magnitude, real part and pointwise complex error. The zooms make peak differences visible. A visually similar magnitude curve alone cannot establish phase or derivative accuracy.}
\clearpage
\subsection{Why an $H^2$ audit matters}
\fig{spatial-pinn-mode-error}{1}{Endpoint Fourier error power weighted by $(1+k^2)^2$ and normalized by reference $H^2$ energy. The unconditioned errors concentrate in the upper half of the retained correction modes. The reference projection floor through mode 48 is negligible in this spatial experiment, so truncation alone cannot explain these errors.}
At $\xi=4$, modes 25--48 account for approximately $93\%$, $84\%$ and $90\%$ of unconditioned $H^2$ error power for packet, sideband and random inputs. Scaling improves this portion of the spectrum without filtering the reported error.
\fig{spatial-pinn-residual-audit}{.82}{Full PDE residual RMS divided by initial RHS RMS on an independent fine spatial grid. One value is shown per fitted slab. These ratios quantify the equation defect, not a certified solution-error bound.}
A small sampled residual must be assessed alongside the independent solution error and spectrum. The reported fits use reference labels only after training; checkpoint selection uses the residual.
\fig{spatial-pinn-phase-curvature}{1}{Phase error is masked where either amplitude is below 5\% of its own peak. Curvature error is unmasked and divided by the reference curvature RMS, not by its full $H^2$ norm. These are supplementary diagnostics; they do not retroactively change the original pass/fail rule.}
For the conditioned packet and sideband fits, curvature-only RMS relative errors remain about $18.0\%$ and $10.6\%$, despite their full $H^2$ errors being below $5\%$. The random curvature error is about $1.53\%$. Thus the next campaign should report a dedicated curvature diagnostic when crest shape or derivative evolution is scientifically important. It should preserve the predeclared primary gate and report additional metrics transparently.

\clearpage
\section{A practical research programme}\label{sec:programme}
\subsection{First ensure that the neural trial space can represent the target}
The old spatial $K=48$ architecture must not be copied unchanged into the temporal $T=16$ task. For a candidate cutoff $K$, compute the best possible Fourier approximation floor
\begin{equation}
 F_{K,2}=\max_{T\in\mathcal T}
 \frac{\norm{(I-P_K)u_{\rm ref}(T)}_{H^2}}{\norm{u_{\rm ref}(T)}_{H^2}}.
 \label{eq:floor}
\end{equation}
It bounds the error of every prediction supported on these modes, regardless of training. At $\epsilon=0.05$, $K=48$ floors through $T=16$ are $42.6\%$, $48.5\%$ and $59.9\%$ for packet, wavetrain and random data. Thus a 5\% neural gate is unattainable in this trial space. Even $K=128$ leaves a roughly 31\% random-profile floor. By $T=4$, all six primary cases instead reach a floor below 5\% using at most $K=16$ among tested cutoffs. Begin on shorter intervals and choose a projection budget strictly below the final error gate.
\IfFileExists{figures/temporal-representation-bandwidth.pdf}{\fig{temporal-representation-bandwidth}{.72}{Representation and physical-band diagnostics for temporal trajectories. Fourier projection floors are measured from reference fields before choosing a neural architecture. A diagnostic threshold $\epsilon|k|\ge0.5$ identifies short modulation scales; it is not a universal validity boundary.}}{}
Narrowband asymptotics require modulation wavenumbers small relative to the carrier. Little wave action outside that regime can dominate $H^2$; numerical convergence alone does not establish physical validity. Report both action and derivative-weighted fractions.

\subsection{What a temporal PINN transfer would require}
Retain the slab construction, replace its linear flow and complete nonlinear residual by \eqref{eq:scaled}, and enforce its explicit model ID. Differentiate in $T$ with automatic differentiation and in $s$ with Fourier multipliers. Compute the density and nonlocal multiplier on a sufficiently padded full spatial grid at every residual time; random isolated space--time points do not by themselves evaluate the nonlocal term. Test Torch and reference RHS parity, plane-wave phase, MI drift and zero-$\epsilon$ agreement before training.

Start at $T\le4$, then $8$, extending only after reference and representation checks pass. Use matched scaled/unscaled fits and a finer independent residual grid. This temporal neural transfer is specified here but has \emph{not} been implemented or trained by the present CPU proof of concept.
\clearpage
\subsection{Test purposes, controls and decisions}
\small
\begin{longtable}{p{.15\linewidth}p{.29\linewidth}p{.28\linewidth}p{.17\linewidth}}
\toprule Test & Scientific purpose & Control and measurement & Decision \\\midrule\endhead
Equation audit & Resolve signs, variables and mean flow & Plane-wave frequency, Fourier symbol, complex MI eigenvalue & Must pass before training \\
Reference refinement & Separate dynamics from discretization & Double grid; halve step independently; $H^0,H^1,H^2$, action, tails & Stop or refine failed intervals \\
Packet box size & Distinguish isolated-packet behaviour from periodic interaction & Double $L,N$ at fixed spacing; matched initial spectrum & Qualify domain claim only through tested time \\
Representation audit & Avoid impossible neural accuracy targets & Compute $F_{K,2}$ using qualified reference & Choose $K$ with floor well below error budget \\
Slab PINN ablation & Test conditioning, including derivative fidelity & Same inputs and budgets; $w_m=1$ vs $(1+k_m^2)^{-1}$; 3 seeds & Extend only when full gates pass \\
Input sensitivity & Quantify finite-time amplification of data differences & Paired bandlimited perturbations at $10^{-4},10^{-3},10^{-2}$; 5 directions & Separate physical sensitivity from model error \\
FNO whole-profile test & Learn a reusable solution operator across data & Entire profiles and all their time windows kept in one split & Compare held-out rollout errors and full costs \\
Later-time extension & Test accumulated error and repeated focusing & No reference resets; evaluation at 4,8,16,32 only when qualified & Report first failed interval; do not infer blow-up \\
\bottomrule
\end{longtable}\normalsize
These are sequential tests, not a promise that every horizon will succeed. The first four are inexpensive CPU prerequisites. The existing spatial PINN ablation is the most evidence-supported neural next step; an operator-learning campaign should follow once reference quality and bandwidth are established.

\subsection{A bounded initial-data ensemble}
A feasible first operator pilot is 16 training, 4 validation and 8 test profiles per family: 48/12/24 trajectories total. These are proposed sample counts, not a statistical power calculation. For example, vary packet peak in $[0.6,0.85]$, width in $[2.5,4]$ and chirp in $[-0.08,0.08]$; vary carrier peak in $[0.6,0.8]$, modulation depth in $[0.02,0.08]$ and mode 3--6; vary saved random Fourier coefficients with bounded peak and bandwidth. Use separate fixed seeds for each split. A new random realization is a new profile, not a new time slice of a training profile.

Do not silently omit hard or unresolved profiles from a purported generalization result. Record numerical-label exclusions and the resulting domain of the dataset. Evaluate different initial phases and translations as symmetry checks, but keep augmented copies of the same base profile in the same split. Start with one temporal formulation and a fixed period; a different $c$ or model ID is a separate experiment.
\clearpage
\subsection{FNO implementation and a fair comparison}
An FNO is intended to learn a map over a family of inputs \cite{fno}; the existing per-inlet PINN is not such an operator. A practical temporal map is
\[
 \mathcal G_\theta:(\operatorname{Re}u_j,\operatorname{Im}u_j,\epsilon,\Delta T)
       \longmapsto u_{j+1}.
\]
Use real and imaginary channels with constant parameter channels, periodic Fourier layers and a residual update around the exact linear flow. Keep normalization fixed from training data. A starting architecture of four Fourier blocks, width 64 and a validated mode count is an engineering choice to tune on validation data, not a scientifically established optimum. The output grid and Fourier blocks must be audited against \eqref{eq:floor}; the block's low-mode count alone does not define the output bandwidth when local layers generate higher modes.

Train with complex-field error and a declared derivative-sensitive loss, for example $L^2$ plus normalized $H^1/H^2$ terms, and short multi-step rollout loss. Include the unweighted-loss baseline to isolate the effect. A physics residual penalty is an optional later ablation; it requires the same full nonlocal evaluation and does not turn residual minimization into an existence proof.

\textbf{Evaluation protocol.} Freeze architecture and checkpoints using validation profiles. Evaluate untouched test profiles under both one-step prediction and autonomous rollout, where each predicted state becomes the next input. Report per-family medians, worst cases, failure counts, peak timing/amplitude errors, complex phase error, $E_0,E_1,E_2$, action drift, spectrum and sensitivity gain. Report a defined valid horizon, such as the first saved station at which any primary gate fails. A low teacher-forced one-step error cannot establish a good autonomous rollout.

Compare the FNO with the numerical solver at \emph{matched accuracy}, including label generation and training. If $C_{\rm setup}$ is reference generation plus training cost, $C_{\rm num}$ is numerical cost per requested trajectory and $C_{\rm infer}$ is model inference cost, the simple amortization estimate is
\[
 n_{\rm break-even}=\frac{C_{\rm setup}}{C_{\rm num}-C_{\rm infer}}
 \quad\text{only if }C_{\rm num}>C_{\rm infer}.
\]
Use measured common-hardware costs and include failed trials. For three profiles, a fast Fourier solver is likely the useful production baseline; no neural speed advantage is established here. A reusable operator becomes scientifically useful if many held-out queries meet derivative and peak requirements at an acceptable total cost. Inverse inference from sparse observations is a plausible later use of PINNs \cite{raissi}, but it would require a new noise/identifiability experiment rather than being claimed from these forward fits.

\subsection{A100 sequence while training is pending}
The supplied cluster handoff pins the existing \emph{spatial} repair revision. Run its smoke configuration first, then the conditioned check and matched unconditioned control using new output directories. Compare three seeds before extending the interval. Use one A100 per task initially; independent profiles or seeds can form an array. Multi-GPU training is unnecessary until a measured memory or throughput bottleneck justifies it. The new temporal reference solver is CPU-based and has a separate \texttt{med} array launcher. No cluster job was submitted as part of this report.

\clearpage
\section{Connection to analytical PDE questions}\label{sec:analysis}
\subsection{A useful distinction that can be proved}
The following elementary fact explains why a smooth-looking finite-dimensional trajectory is weak evidence for continuum global existence.
\begin{proposition}[Mass and finite Fourier systems]
For a smooth periodic solution of \eqref{eq:scaled} with real $c,\epsilon$, $M(T)=\int_0^L|u|^2\,\rad s$ is constant. An exact orthogonal finite Fourier Galerkin truncation of this equation has a global-in-time solution for every finite initial coefficient vector.
\end{proposition}
\begin{proof}
Differentiate $M$ and take the real part of $2\int\overline u u_T$. Periodic integration by parts cancels the skew-adjoint linear terms. The cubic and nonlocal mass pairings are purely imaginary, since $|D_s|(|u|^2)$ is real. The real parts of both derivative nonlinearities are constant multiples of $\int\rho\rho_s$, which vanishes by periodicity. Hence $M'=0$. For a Fourier projection $P_K$, the same pairing holds because $P_K$ is orthogonal and $u_K=P_Ku_K$. The Galerkin vector field is a polynomial finite-dimensional ODE. Conservation keeps its coefficient vector in a bounded mass ball, so local ODE continuation prevents finite-time escape.
\end{proof}
The estimate
\[
 \norm{u_K}_{H^j}\le(1+k_K^2)^{j/2}\norm{u_K}_{L^2}
\]
depends on the cutoff. It gives no uniform continuum derivative bound as $K\to\infty$. The proposition also does not guarantee stability of an explicit time-discretization; the observed time-step artifact illustrates this distinction.

\subsection{What is still needed for an existence or stability statement}
A rigorous continuum claim must specify the exact envelope variable, coefficients, periodic or whole-space domain, regularity class, data size and solution concept. Dysthe well-posedness literature supplies important results for particular formulations and domains \cite{grande,mosincat,heller}; those hypotheses do not automatically match the present one-dimensional periodic potential model. This report does not assert that a matching theorem is absent. It does not supply one.

For a PINN candidate $v$, define the defect $r=v_T-F(v)$ and initial mismatch $v(0)-u_0$. A residual-to-solution estimate would be valuable, but it requires a stability argument for this derivative and nonlocal PDE. \emph{Conditional} on an estimate in a specified norm $X$,
\[
 \frac{\rad}{\rad T}\norm{u-v}_X\le C(T)\norm{u-v}_X+\norm{r}_X,
 \qquad C\ge0,
\]
Gronwall yields
\[
 \norm{u-v}_X(T)\le e^{\int_0^T C(s)\,\rad s}
 \left[\norm{u_0-v(0)}_X+\int_0^T\norm{r(s)}_X\,\rad s\right].
\]
Establishing the hypothesis, including derivative loss and suitable bounds on $u,v$, is the hard analytical step. A sampled residual and a constant borrowed from cubic NLS do not establish it. A sensible collaboration is to use numerical spectra and residuals to identify tractable assumptions, then seek an a posteriori stability result for the chosen formulation.
\clearpage
\subsection{Solitons, recurrent packets and model validity}
A soliton claim needs more than a localized bright curve. A practical first question is whether there is a travelling coherent profile
\[
 u(s,T)=e^{-\ii\Omega T}Q(s-VT)
\]
for the \emph{specified} Dysthe model. Substitute this ansatz, solve the stationary nonlocal profile equation with phase and translation conditions, and refine its residual independently. Compare a propagated candidate after the predicted shift and phase, while also reporting the unaligned error. Perturb it within the same formulation and measure orbital deviation. This is a coherent-wave and stability programme, not a claim of integrability or elastic soliton collisions.

Numerical travelling-wave work exists for spatial Dysthe-type formulations \cite{fedele}; finite-bandwidth and mean-flow effects also matter in temporal wave-group simulations \cite{shemer}. Neither supplies a ready-made exact soliton for every equation labelled Dysthe. The Gaussian input in this report is a packet, not an exact NLS or Dysthe soliton. The NLS sech test verifies only the cubic limiting solver.

The supplied Schober--Islas 2026 preprint studies higher-order weakly viscous extensions and a phase-lag mean-flow damping mechanism \cite{schober}. It is relevant to questions about persistent downshift and dissipative coherent structures, but it changes the model and the conservation structure. No damping or complex mean-flow multiplier is included in the present inviscid baseline. A dissipative comparison would require a separately versioned equation, its own calibration and different conservation diagnostics.

\textbf{Useful analytical discussion questions.}
\begin{enumerate}
\item Which local theory and continuation norm match the selected potential-envelope equation on the periodic domain? Is the later spatial problem analytically preferable for the intended experiment?
\item Can the reference spectra guide a bound on the derivative nonlinearities or a residual-to-solution estimate on a finite interval?
\item Which reported features survive changes of period and narrowband regime, rather than merely a grid refinement?
\item Is a travelling-profile continuation study a more focused first question than a broad claim about generic long-time smoothness?
\end{enumerate}

\subsection{Implementation agreement for a linearized control}
\fig{temporal-mi-check}{.95}{Independent small-perturbation MI check, amplitude $0.8$ and modulation wavenumber $2\pi\cdot5/40$. Twelve trajectories span both $c$ values, both $\epsilon$ values and three grid/step settings. Growth and phase rate agree with the exact retained-model eigenvalue. Applying the wrong $c$ phase formula fails while growth remains identical. This is an implementation check, not a nonlinear stability theorem.}
The perturbation amplitude is $10^{-6}$ and the interval is $T\le4$. Maximum measured rate discrepancies are below $7\cdot10^{-12}$, well below the declared $2\cdot10^{-6}$ absolute tolerance. Such cancellation-limited agreement is not an observed fourth-order convergence rate or a bound on later nonlinear error.

\clearpage
\section{Reproducibility and handoff}\label{sec:repro}
\subsection{What is included}
The editable report source includes all plotted data, vector PDF figures, plotting scripts, the equation audit and references. The separate proof-of-concept archive includes the CPU solver, tests, configurations, exact MI check, result summaries, failed-gate records, checksums, portable plot arrays and cluster instructions. Full reference arrays remain in the local experiment directory; compact arrays in the archive are explicitly visualization-only and must not be used as replacement numerical references.

\begin{center}\small
\begin{tabular}{p{.34\linewidth}p{.57\linewidth}}\toprule
File or directory & Purpose \\\midrule
\texttt{run.py}, \texttt{test\_run.py} & Temporal solver and ten exact/algebraic/convergence tests \\
\texttt{mi\_check.py} & Independent growth and phase control \\
\texttt{primary-results.json} & Six primary cases with explicit per-case provenance \\
\texttt{verification.json} & Runtime, source and array checksums, recorded checks \\
\texttt{results*/summary.json} & Coarse, refined and high-resolution outcomes, including failures \\
\texttt{portable-primary/} & Compact display arrays and full-grid diagnostic summaries \\
\texttt{cluster/README.md} & CPU reference and separate pinned spatial A100 instructions \\
\texttt{experiment-plan.json} & Staged test purposes, metrics and implementation status \\\bottomrule
\end{tabular}
\end{center}
In a prepared Python environment, the basic CPU sequence is:
\begin{verbatim}
python -m pip install -r requirements.txt
python -m unittest -v test_run.py
python run.py --config config-refined.json --output fresh-refined
python run.py --config config-audit.json --output fresh-audit
\end{verbatim}
Use a new output directory when changing source or configuration. The identity check intentionally rejects mismatched cached results. The cluster handoff gives the two-epsilon CPU array and the more expensive time-step stress repair. It uses the supplied \texttt{med} CPU partition; it does not require an A100. Site paths and allocation details must be supplied locally.

\subsection{Existing spatial revisions and evidence status}
The A100 handoff targets the published repair branch \texttt{codex/pinn-h2-rerun} and these revisions:
\begin{center}\small
\begin{tabular}{ll}\toprule
Repository & Revision prefix \\\midrule
\texttt{dysthe-experiments} & \texttt{f53afd9e8538} \\
\texttt{dysthe-learning} & \texttt{bdfe60e571e9} \\
\texttt{dysthe-core} & \texttt{a47a42321475} \\\bottomrule
\end{tabular}
\end{center}
Full commit identifiers and launch commands are in the handoff. The local spatial evidence analysis checked 60 raw-array, summary and selected-checkpoint hashes. This report reanalyses recorded CPU data; it does not claim new PINN training, repaired GPU accuracy or a completed FNO comparison. A job's successful exit is not the same as passing its scientific error gates.

\textbf{For the meeting.} Present the checked temporal focusing examples, the time-step artifact, the finite-mode representation limit and the paired spatial conditioning result. Ask for agreement on the exact target equation and the first finite-time validation window. These are actionable findings even while the repaired cluster campaign is pending.

\clearpage
\small
\begin{thebibliography}{99}
\bibitem{dysthe} K. B. Dysthe, Note on a modification to the nonlinear Schr\"odinger equation for application to deep water waves, \emph{Proc. R. Soc. Lond. A} \textbf{369} (1979), 105--114. User-supplied scan; especially equations (2.6), (2.18)--(2.21), (3.2)--(3.7), (4.1)--(4.3). \url{https://doi.org/10.1098/rspa.1979.0154}.
\bibitem{shemer} L. Shemer and B. Dorfman, Experimental and numerical study of spatial and temporal evolution of nonlinear wave groups, \emph{Nonlinear Processes in Geophysics} \textbf{15} (2008), 931--942. Potential/surface sign discussion on pp. 932--933. \url{https://doi.org/10.5194/npg-15-931-2008}.
\bibitem{trulsen} K. Trulsen, Weakly nonlinear and stochastic properties of ocean wave fields: Application to an extreme wave event, in \emph{Waves in Geophysical Fluids} (2006), Sect. 4.3. \url{https://doi.org/10.1007/978-3-211-69356-8_2}.
\bibitem{fedele} F. Fedele and D. Dutykh, Hamiltonian description and traveling waves of the spatial Dysthe equations, preprint (2011, revised 2012). Used here for the spatial equation and numerical coherent-wave context, not to resolve the temporal sign convention. \url{https://arxiv.org/abs/1110.3605}.
\bibitem{grande} R. Grande, K. M. Kurianski and G. Staffilani, On the nonlinear Dysthe equation, \emph{Nonlinear Analysis} \textbf{207} (2021), 112292. \url{https://arxiv.org/abs/2006.13392}.
\bibitem{mosincat} R. Mosincat, D. Pilod and J.-C. Saut, Global well-posedness and scattering for the Dysthe equation in $L^2(\mathbb R^2)$, \emph{J. Math. Pures Appl.} \textbf{149} (2021), 73--97. Formulation and small-data assumptions must be checked before application. \url{https://arxiv.org/abs/2007.01613}.
\bibitem{heller} G. Heller and C. Shao, Strichartz and multi-linear estimates for the one-dimensional periodic Dysthe equation, preprint (2021). \url{https://arxiv.org/abs/2112.12734}.
\bibitem{raissi} M. Raissi, P. Perdikaris and G. E. Karniadakis, Physics-informed neural networks: A deep learning framework for solving forward and inverse problems involving nonlinear partial differential equations, \emph{J. Comput. Phys.} \textbf{378} (2019), 686--707. \url{https://doi.org/10.1016/j.jcp.2018.10.045}.
\bibitem{causal} S. Wang, S. Sankaran and P. Perdikaris, Respecting causality for training physics-informed neural networks, \emph{Comput. Methods Appl. Mech. Engrg.} \textbf{421} (2024), 116813; preprint 2022. \url{https://arxiv.org/abs/2203.07404}.
\bibitem{failure} A. Krishnapriyan et al., Characterizing possible failure modes in physics-informed neural networks, \emph{NeurIPS} (2021). \url{https://arxiv.org/abs/2109.01050}.
\bibitem{fno} Z. Li et al., Fourier neural operator for parametric partial differential equations, \emph{ICLR} (2021); preprint 2020. \url{https://arxiv.org/abs/2010.08895}.
\bibitem{schober} C. M. Schober and A. Islas, Higher-Order Extensions of Weakly Viscous Dysthe Theory and a Phase-Lag Model for Nonlinear Mean-Flow Damping, user-supplied preprint, version 1, 5 August 2026. \url{https://arxiv.org/abs/2608.04473}.
\end{thebibliography}
\paragraph{How the literature is used.}
Classical PINNs and FNOs motivate the numerical architectures; they do not validate these methods for this exact water-wave equation. Causal and sequence-based PINN work \cite{causal,failure} motivates controlling temporal training difficulty. Our hard slab handoff is not identical to causal loss weighting. No general neural-method speed or accuracy claim is transferred from another PDE. The strongest Dysthe-specific evidence in this report is the explicitly labelled computation and source audit.

\end{document}
